1.Ans: Java is an object oriented programming language . The java is known as the platform independent because , when the java program is compiled the class creates a bytecode or .class , this bytecode is run by the JVM (Java Virtual Machine) on any other system which having the JVM . So , the java programming language is known as the platform independent.

2.Ans: Architecture of the JDK :
	JDK consists of two parts , 1.Tools : here the javac is present (javac is used to compile the code )
				    2.JRE : here all the inbuilt libraries , execution of the program all are present in the jre  

3.Ans: JVM stands for the Java Virtual Machine . It is a storage space which is built to execute the program . 	In java , the JVM is used to execute the bytecode or .class .

4.Ans: Write Once Run Anywhere : Java programming language is having the property called "Write Once Run Anywhere" . This means when the java program is written and compiled once , it creates the .class , this  .class can execute or run in the same system or any systems where the JVM is present .

5.Ans: bytecode in Java :  In java , when we write the program , first we create a class . When we compile this program , it creates .class , this .class is called as the bytecode . This bytecode is generated so that the JVM can understand the program which we are written. Because of this bytecode , the java is known as the platform independent. 

6.Ans: DataTypes in java : DataType is used for the declaration of the Variables . The Data type is used to specify what type of the values can be store in the variable and how much the memory is allocated to that variable . When the declaration of the variable is done by giving the specific data type , only that type of values can be store in that variable , if we try to store the different type of values then it throws an error. 
	There are two types of Data Type : 1. Primitive Data type and 2.Non Primitive Data Type 
	1. Primitive Data Type : These are the keywords . There are total 8 primitive data type . They are,
		byte
		short
		int
		long
		float
		double
		boolean  
		char
	Example: int number; 
  		 byte age;
		 long phoneNumber;
	2. Non Primitive Data type : All the class name is a Data type  , those are the non primitive datatype . The non primitive data type are the user defined data type .
	Example: String name;
		 Map map;
		 Cup cup;

7.Ans: Default values :  The values which is assigned by the jvm to the variables if the values are not assigned by the users is called as the default values . Or if the developer is declaring and initializing the variables , but the developer dosen't want to store the values , he can assign the default values 
Program to print all the default values of all datatype (using the instances) 
class PersonDetails{
	byte age;
	short values;
	int id;
	long mobileNumber;
	float marks;
	double totalMarks;
	boolean isStudent;
	char grade;
	String name;
	public 	PersonDetails()
	{
		System.out.println("Creating the constructor with no-args");
	}
}

class PersonDetailsRunner{
	public static void main(String[] values)
	{
		PersonDetails details = new PersonDetails();
		System.out.println("Default value of age: "+details.age+"Default value of values : "+details.values+"Default values of id is : "+details.id+"default values of mobileNumber : "+details.mobileNumber+"Marks :"+details.marks+"totalMarks : "+details.totalMarks+"isStudent : "+details.isStudent+"grade : "+details.grade+" name :"+details.name);
	}
}

Example :
class PersonDetails{
	public static void main(String[] values){
	byte age = 0;
	short values = 0;
	int id = 0;
	long mobileNumber ; 0L;
	float marks = 0.0f;
	double totalMarks = 0.0;
	boolean isStudent = false;
	char grade = '';
	String name = null
	System.out.println("Default values are "+age);
	System.out.println("Default values are "+values);
	System.out.println("Default values are "+id);
	System.out.println("Default values are "+mobileNumber);
	System.out.println("Default values are "+marks);
	System.out.println("Default values are "+totalMarks);
	System.out.println("Default values are "+isStudent);
	System.out.println("Default values are "+grade);
	System.out.println("Default values are "+name);
	}
	
8.Ans: Variables in Java : Variables are the containers to stores the values .
In variables there are 4 types : Static , non static (instance) , local , parameters
The static variables that are declaried inside the class with the static keyword 
The local variables are the variables which are declaried inside mathods or block 
Example of static variables are 
class Map{
	static String place;
	static double coordibnates;
}
Example of local variable
class Cup{
	public static void main(String[] values)
	{
		double size;
		String color;
	}
}

9.Ans : Static method are the method which executes the block of code . Static method can be called from the class name followed by the . operator and then the method name and then parameters (if any).
Syntax of calling static method is :
						className.methodName(arguments);
						
10.Ans :Program
class Sales{
	public static void main(String[] values)
	{
		double day1 = 4080.0;
		double day2 = 3211.0;
		double day3 = 6500.0;
		double day4 = 3544.0;
		double day5 = 2450.0;
		double day6 = 5570.0;
		double day7 = 9500.0;
		double[] days = {day1 , day2 , day3 , day4 , day5 , day6 , day7};
		double week = day1+day2+day3+day4+day5+day6+day7;
		System.out.println("Total sales of the week is : "+week);
	
		for (int start = 0 ; start<days.length;start++){
			for(int value = 1 ; value<days.length;value++){
				if(days[start]<days[value])
				{
					System.out.println("Highest sales is : "+days[value]);
				}
			}
			
		}
	}
}

 
 